\iflong
\ds{in the first part of the paper we deal with traces and bisimilarity. 
The main reason is to emphasize that the correspondence OGS/\piI is robust, and does not
depend on a specific behavioural relation. Of course another reason is to show that the
trace and bisimilarity semantics for OGS  and pi are the same, but this is not really the
main reason, i think
 }





\finish{the operational correspondence OGS/pi, does offer a new look 
into the behaviour of Milner's encoding. It is something we could stress. 
} 


\fi


It would be interesting to see how to  extend the full abstraction result 
Theorem~\ref{t:fa_CT_piI} 
in \piI  
(about the encoding of pure $\lambda$-terms) 
to the whole  $\lambdarho$-calculus.  The result presented exploits the correspondence
between \WBOGS and \COGS, which does not hold in the whole calculus. Recent
studies on the representation of references in $\pi$-calculus \cite{HirschkoffPS20} may
help.
We also leave for future work a study of the algebraic and behavioural 
 properties of  the equivalence of \complete traces in $\pi$-calculus  introduced in that theorem.

\iflong
\finish{more for future work: 
We also leave for future work a study of the algebraic and behavioural 
 properties of  the equivalence of \complete traces, 
introduced in Section~\ref{s:fapiI}.  It would be also interesting to see how to formalise
the  same idea  of `complete observations' onto coinductive
definition of process equivalence such as   bisimilarity  (or barbed congruence).
}

\fi



\iflong
\ds{we should make it clear that the results above mention both traces and bisimilarity
  but traces is the standard thing in OGS (perhps in other sections we only consider traces)}
  \fi
  
\iflong  
  \subsection{Persistent continuations}
\label{s:linearCon}

If the  $\lambda$-calculus  is extended with 
\emph{control operators},  the linearity on the manipulation of 
the continuation  of a $\lambda$-term does not hold: the continuation  may be dropped or
duplicated, an arbitrary number of times. 
In the OGS and $\pi$-calculus representations in Sections~\ref{s:} and \ref{...}, this is
accommodated as follows. In the OGS model,  
  the continuation entries of environments should become persistent; this affects
  rule OA, wher  the environment   on the right-hand side of the transition should remain
  the same as that on the  left-hand side. 
 In the $\pi$-calculus,  the continuation
names used in the encoding  become persistent 
by adding a replication in front of the inputs 
 $\inp q y$ and  $ \inp r w$, in the definition  of application.   
No modification is required in the encoding of OGS configurations and environments in
Figure~\ref{f:OGS_pi_cbv}. 
We call $\cbvSymbNoLin$ this variant of the encoding 
   $\cbvSymb$.  
%\finish{maybe it is not even necessary to introduce the name   $\cbvSymbNoLin$} 
All results in the previous Sections~\ref{s:..}, in particular Theorem~\ref{t:} and \ref{????}, 
continue to hold   under these modifications. 
\finish{more results to emphasize? more to say?} 

\fi
